 >>?% ======================================================================
% CAPITULO 4 --- O Modelo Classico de Mercado
% Baseado em: Farinelli, S. (2021). arXiv:0910.1671v10, Secao 2
% ======================================================================
\chapter{O Modelo Cl issico de Mercado}

Este cap -tulo apresenta o arcabou SSo matem itico padr o da teoria de
arbitragem em tempo cont -nuo, conforme estabelecido por
Harrison--Kreps--Pliska (1979--1981) e consolidado por Delbaen e
Schachermayer (1994, 2006). Todas as defini SS ues, nota SS ues e teoremas
seguem a Se SS o~2 de Farinelli (2021). O leitor que domina este
cap -tulo possui os pr c-requisitos formais para a Parte~II (GAT).

% ----------------------------------------------------------------------
\section{Espa SSo de Probabilidade Filtrado e Condi SS ues Usuais}
% ----------------------------------------------------------------------

O alicerce de qualquer modelo financeiro em tempo cont -nuo c um
\textbf{espa SSo de probabilidade filtrado} que satisfaz as condi SS ues
t ccnicas m -nimas para que a integra SS o estoc istica funcione.

\begin{defi}[Espa SSo de Probabilidade Filtrado]
Um espa SSo de probabilidade filtrado c uma qu idrupla
$(\Omega, \mathcal{A}, \mathbb{A}, P)$, onde:
\begin{enumerate}[(i)]
 \item $(\Omega, \mathcal{A}, P)$ c um espa SSo de probabilidade
 completo --- toda a teoria de integra SS o de Lebesgue se aplica;
 \item $\mathbb{A} = (\mathcal{A}_t)_{t \in [0, +\infty[}$ c uma
 \textbf{filtra SS o}: uma fam -lia crescente de sub-$\sigma$- ilgebras
 de $\mathcal{A}$, isto c, $\mathcal{A}_s \subseteq \mathcal{A}_t$
 para $0 \leq s \leq t < +\infty$;
 \item A filtra SS o satisfaz as \textbf{condi SS ues usuais} de
 Dellacherie (1972)\footnote{Dellacherie, C. \textit{Capacit cs et
 Processus Stochastiques}. Springer, 1972. ISBN: 978-3540056768.
 \textbf{Resenha:} Obra fundacional que introduziu as condi SS ues
 usuais para filtra SS ues em tempo cont -nuo, indispens iveis para a
 teoria geral dos processos.}:
 \begin{enumerate}
 \item \textbf{Continuidade direita:}
 $\mathcal{A}_t = \bigcap_{s > t} \mathcal{A}_s$
 para todo $t \geq 0$;
 \item \textbf{Completude:} $\mathcal{A}_0$ (e portanto
 toda $\mathcal{A}_t$) cont cm todos os conjuntos
 $P$-negligenci iveis de $\mathcal{A}$.
 \end{enumerate}
\end{enumerate}
\end{defi}

\begin{obs}[Interpreta SS o Intuitiva]
$\mathcal{A}_t$ representa \textbf{toda a informa SS o dispon -vel ao
mercado at c o instante} $t$. A condi SS o de crescimento
($\mathcal{A}_s \subseteq \mathcal{A}_t$) reflete o fato de que
informa SS o nunca c perdida. A continuidade direita c uma
regulariza SS o t ccnica que garante que tempos de parada se comportam
bem. A completude assegura que conjuntos de probabilidade zero n o
afetam a estrutura da informa SS o --- eventos imposs -veis n o carregam
informa SS o.
\end{obs}

O horizonte temporal c $[0, +\infty[$. Isto permite modelar mercados
sem data terminal fixa, o que c conveniente para a GAT, onde
transforma SS ues de gauge envolvem integrais sobre $[0, +\infty[$.

% ----------------------------------------------------------------------
\section{Ativos, Pre SSos e a Conta Caixa}
% ----------------------------------------------------------------------

O mercado consiste em $N$ \textbf{ativos de risco} indexados por
$j = 1, \dots, N$, mais um \textbf{ativo livre de risco} que serve
como \textbf{numer irio} (unidade de conta), indexado por $j = 0$.

\subsection{Semimartingales como Modelo de Pre SSos}

A classe de processos estoc isticos que modela os pre SSos deve ser
suficientemente rica para incluir todos os casos de interesse
(movimento Browniano geom ctrico, processos de salto, difus ues com
saltos) e, simultaneamente, suficientemente restrita para que a
integral estoc istica seja bem definida. A classe correta c a dos
\textbf{semimartingales}.

\begin{defi}[Semimartingale]
Um processo adaptado $X = (X_t)_{t \geq 0}$ c um
\textbf{semimartingale} se admite a decomposi SS o
\begin{equation}\label{eq:semimartingale_decomp}
 X_t = X_0 + M_t + A_t,
\end{equation}
onde $M$ c um martingal local (com $M_0 = 0$) e $A$ c um processo
adaptado c dl g de varia SS o finita (com $A_0 = 0$). A decomposi SS o
n o c onica em geral; torna-se onica se exigirmos que $A$ seja
previs -vel (decomposi SS o can 'nica).
\end{defi}

A classe dos semimartingales c maximal para a integra SS o
estoc istica\footnote{Protter, P. \textit{Stochastic Integration and
Differential Equations}. Springer, 2nd ed., 2005. DOI:
10.1007/978-3-662-10061-5. \textbf{Resenha:} Tratamento matem itico
rigoroso da integra SS o estoc istica, cobrindo semimartingales, integral
de It ', teorema de Girsanov e representa SS o de martingais. O
Teorema~II.4 prova que a integral estoc istica s 3 c bem definida
(como operador cont -nuo) exatamente para integradores
semimartingales. Refer ancia t ccnica padr o para pesquisadores em
finan SSas quantitativas.}: o operador integral $I_X: H \mapsto
\int H\,dX$ s 3 c cont -nuo (na topologia de Emery) quando $X$ c um
semimartingale.

\begin{ex}[Processos de It ']
O exemplo mais relevante para aplica SS ues c o processo de It '
$d$-dimensional:
\begin{equation}\label{eq:ito_process}
 dX_t = \mu_t\,dt + \sigma_t\,dW_t,
\end{equation}
onde $W$ c um movimento Browniano $m$-dimensional, $\mu$ c um
processo adaptado localmente integr ivel e $\sigma$ c um processo
adaptado localmente de quadrado integr ivel. Todo processo de It ' c um
semimartingale (com $M_t = \int_0^t \sigma_u\,dW_u$ e
$A_t = \int_0^t \mu_u\,du$).
\end{ex}

\begin{ex}[Processos de Salto (L cvy)]
Processos de L cvy com saltos --- como o processo de Poisson composto
$X_t = \sum_{k=1}^{N_t} Y_k$, onde $N$ c um processo de Poisson e
$Y_k$ s o os tamanhos dos saltos --- tamb cm s o semimartingales.
Estes modelam \textit{crashes} e eventos extremos. A GAT incorpora
naturalmente processos de salto pois a geometria diferencial n o
distingue entre parte cont -nua e parte de salto.
\end{ex}

\subsection{Pre SSos Nominais e Descontados}

\begin{defi}[Pre SSos Nominais]
O vetor de \textbf{pre SSos nominais} c um semimartingale $N$-dimensional
\begin{equation}\label{eq:prices}
 S: [0, +\infty[ \times \Omega \to \mathbb{R}^N,
 \quad S_t = (S_t^1, \dots, S_t^N),
\end{equation}
adaptado filtra SS o $\mathbb{A}$. Cada componente $S_t^j$ c o pre SSo
de uma unidade do ativo $j$ no instante $t$, expresso em unidades
monet irias (reais, d 3lares, etc.).
\end{defi}

A \textbf{conta caixa} (ou \textit{money market account}) $S^0$ c o
ativo livre de risco que acumula juros taxa curta $r_t^0$:
\begin{equation}\label{eq:cash_account}
 S_t^0 := \exp\left(\int_0^t r_u^0\,du\right),
\end{equation}
onde $r^0 = (r_t^0)_{t \geq 0}$ c um processo adaptado localmente
integr ivel. Tipicamente, $r^0$ c determin -stico ou, no m iximo, um
processo de It '. Note que $S_0^0 = 1$ por constru SS o --- uma unidade
monet iria investida na conta caixa no instante zero.

\begin{defi}[Pre SSos Descontados]
Os \textbf{pre SSos descontados} (em rela SS o ao numer irio $S^0$) s o
definidos por:
\begin{equation}\label{eq:discounted}
 \hat{S}_t^j := \frac{S_t^j}{S_t^0}, \quad j = 1, \dots, N,
\end{equation}
com a conven SS o natural $\hat{S}_t^0 \equiv 1$. O vetor descontado c
denotado $\hat{S}_t = (\hat{S}_t^1, \dots, \hat{S}_t^N)$.
\end{defi}

\begin{obs}[Significado Financeiro]
$\hat{S}_t^j$ expressa o pre SSo do ativo $j$ em \textbf{unidades de
caixa no instante} $t$, e n o em reais nominais. Descontar c o
processo cont -nuo de ``trazer a valor presente'': um real no futuro
vale menos que um real hoje. A conta caixa $S_t^0$ c o fator de
capitaliza SS o. Matematicamente, descontar c uma mudan SSa de numer irio
--- o primeiro exemplo de transforma SS o de gauge que a GAT
generalizar i.
\end{obs}

\subsection{O Numer irio}

\begin{defi}[Numer irio]
Um \textbf{numer irio} c um ativo cujo pre SSo c estritamente positivo
($S_t^{\text{Num}} > 0$ q.c. para todo $t$) e que serve como unidade
de conta. A escolha mais comum c $S^0$ (a conta caixa), mas
qualquer ativo estritamente positivo pode desempenhar este papel.
\end{defi}

A GAT explora exaustivamente a \textbf{arbitrariedade do numer irio}:
as leis de mercado devem ser invariantes sob mudan SSa de unidade de
conta. Esta invari ncia --- conhecida como \textbf{simetria de gauge}
--- ser i o ponto de partida para a constru SS o geom ctrica.

% ----------------------------------------------------------------------
\section{Estrat cgias de Negocia SS o}
% ----------------------------------------------------------------------

\subsection{Previsibilidade}

\begin{defi}[Processo Previs -vel]
Um processo estoc istico $x: [0, +\infty[ \times \Omega \to
\mathbb{R}^N$ c \textbf{previs -vel} se c mensur ivel em rela SS o 
$\sigma$- ilgebra previs -vel $\mathcal{P}$, gerada por todos os
processos cont -nuos esquerda e adaptados.
\end{defi}

A previsibilidade codifica a restri SS o fundamental de causalidade:
\textbf{voc a decide sua posi SS o no instante $t$ com base na
informa SS o dispon -vel estritamente antes de $t$}. Voc a n o pode
prever o futuro; voc a n o pode reagir instantaneamente a um salto.
Se o pre SSo salta em $t$, sua estrat cgia $x_t$ j i estava fixada
antes do salto.

\begin{defi}[Estrat cgia de Negocia SS o]
Uma \textbf{estrat cgia de negocia SS o} (ou \textit{portf 3lio}) c um
processo previs -vel $N$-dimensional
\begin{equation}\label{eq:strategy}
 x: [0, +\infty[ \times \Omega \to \mathbb{R}^N,
 \quad x_t = (x_t^1, \dots, x_t^N),
\end{equation}
onde $x_t^j$ c a \textbf{quantidade} (n omero de unidades) do ativo
$j$ mantida no instante $t$. A quantidade $x_t^0$ (posi SS o na conta
caixa) c determinada pela condi SS o de autofinanciamento.
\end{defi}

O \textbf{valor da carteira} (ou \textit{wealth process}) associado 
estrat cgia $x$ c:
\begin{equation}\label{eq:wealth}
 V_t := V_t^x := \sum_{j=0}^N x_t^j S_t^j = x_t^0 S_t^0 + x_t \cdot S_t.
\end{equation}

\subsection{Admissibilidade}

Nem toda estrat cgia previs -vel c economicamente razo ivel. A condi SS o
de \textbf{admissibilidade} exclui estrat cgias que realizam lucros
infinitos com capital finito --- as chamadas ``estrat cgias de
dobrar'' (\textit{doubling strategies}).

\begin{defi}[Estrat cgia Admiss -vel]
Uma estrat cgia $x$ c dita \textbf{$a$-admiss -vel} (para $a \in
\mathbb{R}$) se o processo de valor descontado c limitado
inferiormente:
\begin{equation}\label{eq:admissible}
 \hat{V}_t^x = \frac{V_t^x}{S_t^0} \geq -a, \quad \forall t \geq 0,
 \quad P\text{-q.c.}
\end{equation}
Uma estrat cgia c \textbf{admiss -vel} se c $a$-admiss -vel para algum
$a \in \mathbb{R}_+$. Denotamos por $\mathcal{A}$ o conjunto de todas
as estrat cgias admiss -veis.
\end{defi}

\begin{obs}[Por que limitar a perda?]
Em tempo cont -nuo, sem a limita SS o inferior, c poss -vel construir uma
estrat cgia de ``aposta de Martingale'' que dobra a aposta ap 3s cada
perda e gera lucro certo com probabilidade~1 --- mas requer cr cdito
ilimitado. A cota inferior $a$ representa uma \textbf{linha de cr cdito
finita} e exclui estas patologias\footnote{Harrison, J.M. e Pliska,
S.R. \textit{Martingales and Stochastic Integrals in the Theory of
Continuous Trading}. Stochastic Processes and their Applications,
11(3):215--260, 1981. DOI: 10.1016/0304-4149(81)90026-0.
\textbf{Resenha:} O artigo que estendeu o FTAP de tempo discreto
(Harrison--Kreps 1979) para tempo cont -nuo. Formalizou o conceito de
estrat cgia admiss -vel e provou que, sob a exist ancia de uma medida
martingal equivalente, o valor descontado de qualquer estrat cgia
admiss -vel autofinanciada c um martingal local limitado
inferiormente --- portanto, um supermartingal.}.
\end{obs}

\subsection{Autofinanciamento: It ' e Stratonovich}

A condi SS o de \textbf{autofinanciamento} afirma que varia SS ues no valor da
carteira decorrem exclusivamente de varia SS ues nos pre SSos dos ativos ---
voc a n o injeta nem retira capital.

\begin{defi}[Estrat cgia Autofinanciada --- Formula SS o de It ']
Uma estrat cgia $x$ c \textbf{autofinanciada} (na formula SS o de It ') se
\begin{equation}\label{eq:self_financing_ito}
 V_t = V_0 + \int_0^t x_u \cdot dS_u
 = V_0 + \sum_{j=0}^N \int_0^t x_u^j\,dS_u^j,
\end{equation}
onde a integral c a \textbf{integral estoc istica de It '}. Em nota SS o
diferencial:
\begin{equation}\label{eq:self_financing_diff}
 dV_t = x_t \cdot dS_t = \sum_{j=0}^N x_t^j\,dS_t^j.
\end{equation}
\end{defi}

A integral de It ' avalia o integrando no \textit{in -cio} do intervalo
de integra SS o. Isto a torna n o-antecipativa (o valor da integral no
futuro n o pode ser conhecido hoje) --- uma propriedade essencial para
finan SSas. Contudo, a regra da cadeia usual do c ilculo falha para a
integral de It ', sendo substitu -da pelo \textbf{Lema de It '}, que
adiciona um termo de corre SS o de segunda ordem (o termo de It '):
\begin{equation}\label{eq:ito_lemma}
 df(X_t) = f'(X_t)\,dX_t + \frac{1}{2} f''(X_t)\,d\langle X, X\rangle_t.
\end{equation}

\begin{defi}[Estrat cgia Autofinanciada --- Formula SS o de Stratonovich]
A mesma condi SS o pode ser expressa via \textbf{integral de
Stratonovich}:
\begin{equation}\label{eq:self_financing_stratonovich}
 V_t = V_0 + \int_0^t x_u \circ dS_u
 = V_0 + \sum_{j=0}^N \int_0^t x_u^j \circ dS_u^j.
\end{equation}
A nota SS o ``$\circ$'' indica a integral de Stratonovich, definida pela
rela SS o com a integral de It ':
\begin{equation}\label{eq:ito_stratonovich_relation}
 \int_0^t x_u \circ dS_u
 := \int_0^t x_u \cdot dS_u
 + \frac{1}{2}\langle x, S\rangle_t,
\end{equation}
onde $\langle x, S\rangle_t$ c o processo de \textbf{covaria SS o
quadr itica} entre $x$ e $S$, i.e., a parte cont -nua (previs -vel) da
varia SS o conjunta.
\end{defi}

\begin{obs}[Por que duas formula SS ues?]
A distin SS o entre It ' e Stratonovich n o c um preciosismo t ccnico ---
 c \textbf{fundamental para a GAT}. A integral de Stratonovich satisfaz
a regra da cadeia usual do c ilculo:
\begin{equation}\label{eq:stratonovich_chain}
 df(X_t) = f'(X_t) \circ dX_t,
\end{equation}
exatamente como no c ilculo de Newton--Leibniz. Isto significa que
\textbf{a geometria diferencial (que depende crucialmente da regra da
cadeia) requer Stratonovich}. A GAT formula todas as constru SS ues
geom ctricas (conex ues, curvatura, transporte paralelo) usando
Stratonovich. A integral de It ' permanece a ferramenta correta para
martingais e para a teoria de apre SSamento, mas c Stratonovich que
permite ``colar'' os peda SSos geom ctricos.
\end{obs}

A rela SS o~\eqref{eq:ito_stratonovich_relation} c a ponte entre os
dois mundos: o mundo probabil -stico (It ') e o mundo geom ctrico
(Stratonovich). A covari ncia quadr itica $\langle x, S\rangle_t$ c o
\textbf{termo de conex o} que corrige a ``n o-regularidade'' da
integral de It '.

% ----------------------------------------------------------------------
\section{Arbitragem Formalizada}
% ----------------------------------------------------------------------

Com as defini SS ues anteriores, podemos dar uma defini SS o matematicamente
precisa de arbitragem.

\begin{defi}[Arbitragem]
Uma estrat cgia autofinanciada admiss -vel $x \in \mathcal{A}$ c uma
\textbf{oportunidade de arbitragem} se existe $T > 0$ tal que uma das
duas condi SS ues seguintes se verifica:
\begin{enumerate}[(i)]
 \item $P[V_0^x < 0] = 1$ \;\; e \;\; $P[V_T^x \geq 0] = 1$
 (voc a come SSa com d -vida certa e termina com patrim 'nio n o
 negativo certo);
 \item $P[V_0^x \leq 0] = 1$, \;\; $P[V_T^x \geq 0] = 1$ \;\; e
 \;\; $P[V_T^x > 0] > 0$
 (voc a come SSa sem capital, nunca perde, e tem probabilidade
 positiva de lucro).
\end{enumerate}
\end{defi}

\begin{obs}[Interpreta SS o Financeira]
A condi SS o~(i) descreve a situa SS o em que o mercado literalmente
\textit{paga} voc a para entrar numa estrat cgia que nunca dar i preju -zo
(arbitragem de tipo~I). A condi SS o~(ii) descreve a situa SS o cl issica:
voc a n o investe nada, n o corre risco de perda, mas pode ganhar algo
(arbitragem de tipo~II). Em tempo discreto com espa SSo de estados
finito, a condi SS o~(i) implica a~(ii), mas em tempo cont -nuo as duas
s o independentes.
\end{obs}

\begin{ex}[Arbitragem por Mudan SSa de Numer irio]
Considere dois ativos $S^1$ e $S^2$ com pre SSos estritamente positivos.
Defina $\hat{S}_t := S_t^1 / S_t^2$ (pre SSo de $S^1$ em unidades de
$S^2$). Se $\hat{S}_0 \neq \hat{S}_T$ com probabilidade~1 para algum
$T$, e a dire SS o da varia SS o c conhecida, existe arbitragem:
compre o ativo mais barato e venda o mais caro.
\end{ex}

% ----------------------------------------------------------------------
\section{Condi SS o NA e suas Limita SS ues em Tempo Cont -nuo}
% ----------------------------------------------------------------------

A condi SS o mais b isica de efici ancia de mercado c a \textbf{aus ancia
de arbitragem}:

\begin{defi}[Condi SS o NA --- \textit{No Arbitrage}]
O mercado satisfaz \textbf{NA} se n o existe nenhuma estrat cgia
admiss -vel $x \in \mathcal{A}$ que seja uma oportunidade de
arbitragem.
\end{defi}

Em tempo discreto com espa SSo de probabilidade finito, a condi SS o NA c
equivalente exist ancia de uma medida martingal equivalente
(Harrison--Kreps, 1979)\footnote{Harrison, J.M. e Kreps, D.M.
\textit{Martingales and Arbitrage in Multiperiod Securities Markets}.
Journal of Economic Theory, 20(3):381--408, 1979. DOI:
10.1016/0022-0531(79)90043-7. \textbf{Resenha:} O artigo seminal que
provou pela primeira vez a equival ancia entre aus ancia de arbitragem e
exist ancia de medida martingal equivalente, no contexto de mercados
multi-per -odo com espa SSo de estados finito. Introduziu o conceito de
``pre SSo de estado'' (\textit{state price}) e demonstrou que a
exist ancia de um vetor de pre SSos de estado estritamente positivo c
equivalente condi SS o NA.}.

\textbf{Em tempo cont -nuo, a condi SS o NA c insuficiente.} Existem
mercados que satisfazem NA mas nos quais c poss -vel construir
estrat cgias que ``aproximam'' uma arbitragem --- o risco tende a zero,
mas o lucro esperado permanece positivo. O exemplo can 'nico c a
\textbf{estrat cgia de venda a descoberto com horizonte exponencial}:
venda $e^{-t}$ unidades do ativo no instante $t$; o investimento
inicial converge para zero, o risco converge para zero, mas o payoff
terminal permanece estritamente positivo.

\begin{ex}[NA c Insuficiente --- Esbo SSo]
Considere um ativo cujo pre SSo segue um movimento Browniano geom ctrico
$S_t = \exp(W_t - t/2)$, onde $W$ c um $P$-Browniano. A estrat cgia
$x_t := \mathbbm{1}_{S_t > K}$ para algum $K > 0$ n o c previs -vel.
Toda estrat cgia previs -vel que ``tenta'' capturar esta oportunidade
falha em ser admiss -vel ou gera perda ilimitada com probabilidade
positiva. Contudo, c poss -vel construir uma \textbf{sequ ancia} de
estrat cgias admiss -veis $(x^{(n)})_{n \in \mathbb{N}}$ cujo risco
tende a zero e cujo payoff tende a algo n o-negativo com probabilidade
positiva de ser estritamente positivo. O limite desta sequ ancia n o c
uma estrat cgia admiss -vel --- mas a \textit{exist ancia da sequ ancia} j i
viola a intui SS o econ 'mica de ``mercado eficiente''.
\end{ex}

Para fechar esta brecha, precisamos de uma condi SS o \textbf{topol 3gica}
mais forte.

% ----------------------------------------------------------------------
\section{Condi SS o NFLVR}
% ----------------------------------------------------------------------

Delbaen e Schachermayer (1994)\footnote{Delbaen, F. e Schachermayer, W.
\textit{A General Version of the Fundamental Theorem of Asset Pricing}.
Mathematische Annalen, 300(1):463--520, 1994. DOI: 10.1007/BF01450498.
\textbf{Resenha:} O artigo que estabeleceu o FTAP para tempo cont -nuo
com processos de pre SSo semimartingales. Introduziu a condi SS o NFLVR e
provou sua equival ancia exist ancia de uma medida martingal
equivalente $\sigma$-martingal. Generaliza todos os resultados
anteriores (Harrison--Pliska 1981, Dalang--Morton--Willinger 1990) e c
considerado um dos resultados mais profundos da matem itica financeira
moderna.} introduziram a condi SS o \textbf{NFLVR} (\textit{No Free Lunch
with Vanishing Risk}), que exclui n o apenas arbitragens puras, mas
tamb cm \textbf{aproxima SS ues} de arbitragens.

\subsection{Constru SS o dos Cones}

Considere o conjunto de todos os payoffs terminais descontados
ating -veis por estrat cgias admiss -veis com investimento inicial zero
(ou negativo):

\begin{defi}[Cone $K_0$]
\begin{equation}\label{eq:K0}
 K_0 := \left\{
 \int_0^\infty x_u \cdot d\hat{S}_u \;\middle|\;
 x \in \mathcal{A},\; V_0^x \leq 0
 \right\}.
\end{equation}
$K_0$ c o conjunto de \textbf{payoffs terminais ating -veis sem
investimento inicial positivo}. Cada elemento de $K_0$ c uma vari ivel
aleat 3ria representando o valor terminal (descontado) de uma
estrat cgia que n o exigiu capital inicial.
\end{defi}

\begin{defi}[Cone $C_0$ e Cone $C$]
Definimos:
\begin{align}
 C_0 &:= K_0 - L_+^0(\Omega, \mathcal{A}_\infty, P)
 \label{eq:C0} \\
 C &:= C_0 \cap L^\infty(\Omega, \mathcal{A}_\infty, P)
 \label{eq:C}
\end{align}
onde $L_+^0$ c o cone das vari iveis aleat 3rias n o-negativas
(possivelmente n o limitadas) e $L^\infty$ c o espa SSo das vari iveis
aleat 3rias essencialmente limitadas.
\end{defi}

\begin{obs}[Interpreta SS o dos Cones]
\begin{itemize}
 \item $K_0$: payoffs que voc a pode atingir com capital inicial zero
 (ou negativo). %% o ``card ipio'' de resultados financeiros
 dispon -veis sem investimento.
 \item $C_0$: $K_0$ menos vari iveis n o-negativas. Subtrair algo de
 $L_+^0$ significa que voc a pode \textbf{descartar} dinheiro
 (domin ncia): se voc a pode obter $X$ com capital zero, voc a
 pode obter qualquer $Y \leq X$ (simplesmente jogando fora a
 diferen SSa). Esta c a propriedade de \textbf{free disposal}.
 \item $C$: a restri SS o a $L^\infty$ (payoffs limitados). A raz o
 t ccnica c que $L^\infty$ c o dual de $L^1$, e o teorema de
 separa SS o de Hahn--Banach (essencial para provar a exist ancia
 da medida martingal) requer o espa SSo $L^\infty$ com sua
 topologia natural.
\end{itemize}
\end{obs}

\subsection{Defini SS o Topol 3gica}

\begin{defi}[Condi SS o NFLVR --- Delbaen--Schachermayer 1994]
O mercado satisfaz \textbf{NFLVR} (\textit{No Free Lunch with
Vanishing Risk}) se
\begin{equation}\label{eq:NFLVR}
 \overline{C} \cap L_+^\infty(\Omega, \mathcal{A}_\infty, P) = \{0\},
\end{equation}
onde $\overline{C}$ denota o \textbf{fecho de $C$ na topologia
forte} (norma) de $L^\infty$.
\end{defi}

\begin{obs}[Tradu SS o para Portugu as Claro]
A condi SS o NFLVR diz: \textit{N o existe sequ ancia de payoffs
ating -veis (limitados, com investimento inicial zero ou negativo) que
convirja uniformemente para um payoff n o-negativo que seja
estritamente positivo com probabilidade positiva.} Em outras palavras:
voc a n o pode ``aproximar'' uma arbitragem --- se o risco de uma
sequ ancia de estrat cgias tende a zero (converg ancia em
$L^\infty$), o payoff limite tamb cm deve ser exatamente zero.
\end{obs}

A condi SS o NFLVR decomp ue-se em duas condi SS ues mais simples que s o
 oteis nas demonstra SS ues:

\begin{prop}[Decomposi SS o da NFLVR]
NFLVR c equivalente conjun SS o de:
\begin{enumerate}[(i)]
 \item \textbf{NA} (aus ancia de arbitragem): $C \cap L_+^\infty = \{0\}$;
 \item \textbf{NFL} (\textit{No Free Lunch}):
 $\overline{C}_{\text{frac}} \cap L_+^\infty = \{0\}$, onde o
 fecho fraco estrito (relativo topologia fraca-$\ast$ de
 $L^\infty$) exclui limites de sequ ancias limitadas.
\end{enumerate}
\end{prop}

% ----------------------------------------------------------------------
\section{FTAP em Tempo Cont -nuo}
% ----------------------------------------------------------------------

O \textbf{Primeiro Teorema Fundamental do Apre SSamento de Ativos}
(FTAP) estabelece a equival ancia entre a condi SS o NFLVR e a exist ancia
de uma medida de probabilidade equivalente sob a qual os pre SSos
descontados s o martingais (em sentido generalizado).

\begin{defi}[Medida Martingal Equivalente]
Uma medida de probabilidade $Q$ em $(\Omega, \mathcal{A}_\infty)$ c
uma \textbf{medida martingal equivalente} (EMM --- \textit{Equivalent
Martingale Measure}) se:
\begin{enumerate}[(i)]
 \item $Q \sim P$ (equival ancia: $Q$ e $P$ t am os mesmos conjuntos
 de medida nula);
 \item O processo de pre SSos descontados $\hat{S}$ c um
 $\sigma$-\textbf{martingal} sob $Q$: existe uma sequ ancia
 crescente de tempos de parada $(\tau_n)_{n \in \mathbb{N}}$
 com $\tau_n \uparrow +\infty$ $Q$-q.c. tal que, para cada $j$,
 o processo parado $\hat{S}^{j, \tau_n}$ c um $Q$-martingal.
\end{enumerate}
\end{defi}

\begin{obs}[$\sigma$-martingal vs. Martingal Local]
A no SS o de $\sigma$-martingal c uma generaliza SS o t ccnica de
martingal local necess iria para cobrir todos os semimartingales. Para
processos localmente limitados (i.e., processos de It ' sem saltos
explosivos), $\sigma$-martingal coincide com martingal local. Para
processos com saltos, a distin SS o c relevante. O leitor interessado
nos detalhes t ccnicos deve consultar Delbaen--Schachermayer (2006,
Cap.~9)\footnote{Delbaen, F. e Schachermayer, W. \textit{The
Mathematics of Arbitrage}. Springer Finance, 2006. DOI:
10.1007/978-3-540-31299-4. \textbf{Resenha:} Obra de refer ancia
definitiva sobre os fundamentos matem iticos da arbitragem. Cobre desde
o modelo discreto (Dalang--Morton--Willinger) at c o caso cont -nuo com
semimartingales gerais. O Cap -tulo~9 cont cm a prova completa do FTAP
via teorema de separa SS o de Kreps--Yan. Indispens ivel para qualquer
estudo s crio de finan SSas quantitativas.}.
\end{obs}

\begin{teo}[FTAP --- Delbaen--Schachermayer 1994]
\label{teo:FTAP_DS}
Para um mercado com pre SSos modelados por um semimartingale $S$
localmente limitado, as seguintes condi SS ues s o equivalentes:
\begin{enumerate}[(i)]
 \item O mercado satisfaz \textbf{NFLVR};
 \item Existe uma \textbf{medida martingal equivalente} $Q$ (EMM)
 tal que $\hat{S}$ c um $\sigma$-martingal sob $Q$.
\end{enumerate}
Al cm disso, o conjunto de EMMs com densidade limitada c denso (na
topologia da varia SS o total) no conjunto de todas as EMMs.
\end{teo}

\begin{proof}[Esbo SSo da Prova --- Ideias Centrais]
A implica SS o (ii)~$\Rightarrow$~(i) c a mais f icil e segue do
Teorema de Parada Opcional para $\sigma$-martingais e do Lema de
Fatou.

A implica SS o (i)~$\Rightarrow$~(ii) c profunda e requer:
\begin{enumerate}
 \item NFLVR garante que $C$ c fechado em $L^\infty$ (Lema de
 Kreps--Yan);
 \item Pelo Teorema de Separa SS o de Hahn--Banach, existe um
 funcional linear cont -nuo (um elemento de $(L^\infty)^*$) que
 separa $C$ de $L_+^\infty \setminus \{0\}$;
 \item A representa SS o de Yosida--Hewitt decomp ue este funcional em
 parte absolutamente cont -nua (densidade $dQ/dP$) e parte
 singular;
 \item A condi SS o NA elimina a parte singular, produzindo uma
 densidade estritamente positiva para $Q$;
 \item A condi SS o de $\sigma$-martingal segue da propriedade de
 martingal local da integral estoc istica sob $Q$.
\end{enumerate}
A prova completa ocupa aproximadamente 40 p iginas no artigo original.
\end{proof}

\begin{coro}[FTAP para Processos de It ' Localmente Limitados]
Se $S$ c um processo de It ' localmente limitado (sem saltos, drift e
difus o localmente limitados), NFLVR c equivalente exist ancia de uma
EMM sob a qual $\hat{S}$ c um martingal local.
\end{coro}

% ----------------------------------------------------------------------
\section{Pricing Kernel / State Price Deflator}
% ----------------------------------------------------------------------

Uma formula SS o equivalente --- e particularmente otil para a GAT ---
utiliza o conceito de \textbf{deflator de pre SSos de estado}
(\textit{state price deflator}) ou \textbf{pricing kernel}.

\begin{defi}[State Price Deflator]
Um processo estoc istico estritamente positivo $\beta = (\beta_t)_{t
\geq 0}$ com $\beta_0 = 1$ c um \textbf{state price deflator} (SPD)
se o produto $\beta_t S_t$ c um $P$-martingal local para todo ativo
(incluindo a conta caixa):
\begin{equation}\label{eq:SPD}
 \beta_t S_t^j \text{ c um } P\text{-martingal local}, \quad
 j = 0, 1, \dots, N.
\end{equation}
Equivalentemente, $\beta$ c um \textbf{pricing kernel}.
\end{defi}

\begin{obs}[Rela SS o com a EMM]
Se $Q$ c uma EMM com processo de densidade $Z_t :=
\mathbb{E}_P[dQ/dP \mid \mathcal{A}_t]$, ent o
\begin{equation}\label{eq:SPD_from_EMM}
 \beta_t := \frac{Z_t}{S_t^0}
\end{equation}
 c um state price deflator. Reciprocamente, se $\beta$ c um SPD,
ent o $Q$ definida por $dQ/dP|_{\mathcal{A}_t} := \beta_t S_t^0$ c uma
EMM. A correspond ancia c bijetiva.
\end{obs}

\begin{teo}[Equival ancia NFLVR $\Leftrightarrow$ Exist ancia de SPD]
\label{teo:SPD_equivalence}
Para um mercado com pre SSos modelados por semimartingales localmente
limitados, as seguintes condi SS ues s o equivalentes:
\begin{enumerate}[(i)]
 \item O mercado satisfaz \textbf{NFLVR};
 \item Existe um \textbf{state price deflator} $\beta$ (processo
 estritamente positivo, $\beta_0 = 1$) tal que
 $\beta S^j$ c um $\sigma$-martingal para $j = 0, \dots, N$.
\end{enumerate}
\end{teo}

\begin{proof}
Seja $Q$ uma EMM com processo de densidade $Z$. Defina
$\beta_t := Z_t / S_t^0$. Como $S_t^0$ c determin -stico (ou quase), e
$Z$ c um $P$-martingal, $\beta$ c um $P$-semimartingale. Para cada
$j$:
\[
\beta_t S_t^j = Z_t \frac{S_t^j}{S_t^0} = Z_t \hat{S}_t^j,
\]
que c um $\sigma$-martingal sob $P$ se, e somente se, $\hat{S}_t^j$ c
um $\sigma$-martingal sob $Q$ (pelo Teorema de Girsanov
generalizado). A rec -proca segue revertendo a constru SS o.
\end{proof}

O SPD desempenha um papel central na GAT porque sua derivada
estoc istica (no sentido de Stratonovich) est i diretamente relacionada
 \textbf{conex o de mercado}. Como veremos no Cap -tulo~7, a taxa
curta de cada carteira satisfaz:
\begin{equation}\label{eq:short_rate_SPD}
 r_t^x = -\mathcal{D} \log(\beta_t D_t^x),
\end{equation}
onde $\mathcal{D}$ c a derivada m cdia de Nelson (derivada de
Stratonovich).

% ----------------------------------------------------------------------
\section{Cashflows e Intensidades}
% ----------------------------------------------------------------------

A Se SS o~2.2 de Farinelli (2021) introduz a no SS o de \textbf{cashflows}
e \textbf{intensidades}, que s o fundamentais para a constru SS o dos
\textbf{gauges} no Cap -tulo~5.

\begin{defi}[Cashflow]
Um \textbf{cashflow} c uma fam -lia de vari iveis aleat 3rias
$(C_t)_{t \geq 0}$ onde $C_t$ representa o pagamento (dividendo,
cupom, principal) recebido no instante $t$. Um cashflow c dito
\textbf{admiss -vel} se c um processo adaptado c dl g com
$\mathbb{E}[\int_0^\infty |C_t|\,dt] < \infty$.
\end{defi}

Em tempo cont -nuo, c mais natural trabalhar com a \textbf{intensidade}
(ou densidade) de cashflow, an iloga ``velocidade instant nea'' de
pagamentos.

\begin{defi}[Intensidade de Cashflow]
Uma \textbf{intensidade de cashflow} c um processo estoc istico
$\pi = (\pi_t)_{t \geq 0}$ adaptado e localmente integr ivel. O
cashflow acumulado at c o instante $T$ c:
\begin{equation}\label{eq:cashflow_intensity}
 C_T^\pi := \int_0^T \pi_u\,du.
\end{equation}
Em particular, o pagamento infinitesimal no intervalo $[t, t+dt[$ c
$\pi_t\,dt$.
\end{defi}

\begin{ex}[T -tulo de Cupom Zero (\textit{Zero-Coupon Bond})]
Um t -tulo que paga $1$ unidade monet iria no vencimento $T$ tem
intensidade de cashflow concentrada em $T$:
\begin{equation}\label{eq:zero_coupon_intensity}
 \pi_t = \delta_T(t),
\end{equation}
onde $\delta_T$ c a distribui SS o Delta de Dirac na data $T$
(generalizada). Formalmente, $\int_0^\infty f(t)\,\delta_T(t)\,dt =
f(T)$ para toda fun SS o $f$ cont -nua. O uso de distribui SS ues c
justificado no contexto de funcionais lineares cont -nuos sobre espa SSos
de fun SS ues teste; na GAT, as integrais de gauge envolvem convolu SS ues
que regularizam naturalmente estas distribui SS ues.
\end{ex}

\begin{ex}[A SS o com Dividendos Cont -nuos]
Uma a SS o que paga dividendos a uma taxa cont -nua $d_t$ (propor SS o do
pre SSo) tem intensidade:
\begin{equation}\label{eq:dividend_intensity}
 \pi_t = d_t S_t.
\end{equation}
Se os dividendos s o reinvestidos, a intensidade c nula e o pre SSo
$S_t$ j i incorpora o retorno total.
\end{ex}

\begin{ex}[T -tulo com Cupom]
Um t -tulo com vencimento $T$ que paga cupons anuais de valor $c$ e
principal $1$ no vencimento tem intensidade:
\begin{equation}\label{eq:coupon_bond_intensity}
 \pi_t = \sum_{k=1}^{\lfloor T \rfloor} c\,\delta_{k}(t)
 + (1 + c)\,\delta_T(t).
\end{equation}
\end{ex}

\begin{ex}[Perpetuidade]
Uma perpetuidade que paga $1$ unidade por ano indefinidamente:
\begin{equation}\label{eq:perpetuity_intensity}
 \pi_t = 1, \quad \forall t \geq 0.
\end{equation}
Este c o caso mais simples e serve como exemplo can 'nico na
constru SS o de gauges.
\end{ex}

% ----------------------------------------------------------------------
\section{Valor Presente sob NFLVR}
% ----------------------------------------------------------------------

Sob a condi SS o NFLVR (e, portanto, na presen SSa de um pricing kernel
$\beta$), o valor presente de qualquer cashflow pode ser expresso
de forma onica.

\begin{teo}[F 3rmula de Apre SSamento com Pricing Kernel]
\label{teo:pricing_kernel_formula}
Suponha que o mercado satisfaz NFLVR e seja $\beta$ um state price
deflator. O \textbf{valor presente} (pre SSo de n o-arbitragem) no
instante $t$ de um cashflow com intensidade $\pi$ c:
\begin{equation}\label{eq:present_value}
 V_t^\pi = \frac{1}{\beta_t}\,
 \mathbb{E}_P\!\left[
 \int_t^\infty \beta_u\,\pi_u\,du \;\middle|\;
 \mathcal{A}_t
 \right].
\end{equation}
Em particular, o valor presente de um fluxo unit irio perp ctuo
($\pi_t = 1$) c:
\begin{equation}\label{eq:perpetuity_value}
 D_t := V_t^{\pi=1} = \frac{1}{\beta_t}\,
 \mathbb{E}_P\!\left[
 \int_t^\infty \beta_u\,du \;\middle|\;
 \mathcal{A}_t
 \right].
\end{equation}
\end{teo}

\begin{proof}
Pela defini SS o de $\beta$, o processo $\beta_t S_t^j$ c um
$P$-martingal local para cada ativo $j$. O valor presente de um
cashflow c, por aus ancia de arbitragem, o custo de replic i-lo com uma
carteira autofinanciada. Pelo Teorema de Representa SS o de Martingais
(quando aplic ivel), existe uma estrat cgia $x$ tal que
$V_t^x = V_t^\pi$. Aplicando a defini SS o de SPD:
\[
\beta_t V_t^\pi = \mathbb{E}_P\!\left[
 \int_t^\infty \beta_u\,\pi_u\,du \;\middle|\; \mathcal{A}_t
\right],
\]
j i que $\beta_t V_t^\pi$ deve ser um martingal sob $P$ ( c o valor
descontado pelo SPD de uma carteira autofinanciada). Dividindo por
$\beta_t$ obt cm-se~\eqref{eq:present_value}.
\end{proof}

\begin{obs}[Interpreta SS o Econ 'mica]
A f 3rmula~\eqref{eq:present_value} generaliza o desconto tradicional
($e^{-r(T-t)}$) para um ambiente estoc istico com pricing kernel
vari ivel no tempo e nos estados da natureza. O pricing kernel $\beta$
atua como uma ``taxa de desconto estoc istica instant nea'' que
incorpora tanto o valor temporal do dinheiro quanto o pre SSo de mercado
do risco. Em particular, se $\beta_t = e^{-rt}$ (determin -stico),
recupera-se a f 3rmula cl issica de desconto.
\end{obs}

A quantidade $D_t$ definida em~\eqref{eq:perpetuity_value} --- o valor
presente de uma perpetuidade unit iria --- ser i o \textbf{deflator} no
sentido da GAT. Ela mede, em unidades monet irias de hoje, quanto vale
o direito de receber 1 unidade monet iria por ano para sempre. Este c
o instrumento financeiro mais fundamental, a partir do qual todos os
outros podem ser constru -dos via transforma SS ues de gauge.

\begin{prop}[Propriedade de Martingal do Valor Descontado pelo SPD]
\label{prop:SPD_martingale}
Para qualquer cashflow admiss -vel com intensidade $\pi$, o processo
\begin{equation}\label{eq:SPD_discounted_value}
 M_t^\pi := \beta_t V_t^\pi + \int_0^t \beta_u \pi_u\,du
\end{equation}
 c um $P$-martingal (local). Em particular, para uma perpetuidade
($\pi_t = 1$):
\begin{equation}\label{eq:SPD_martingale_perpetuity}
 M_t := \beta_t D_t + \int_0^t \beta_u\,du
\end{equation}
 c um $P$-martingal, donde se obt cm a din mica:
\begin{equation}\label{eq:D_dynamics}
 dD_t = (r_t^0 D_t - 1)\,dt + dM_t^\beta,
\end{equation}
onde $M^\beta$ c a parte martingal (sob $P$) do processo $\beta D$.
\end{prop}

\begin{proof}
Pela defini SS o de $V_t^\pi$ em~\eqref{eq:present_value}:
\[
\beta_t V_t^\pi = \mathbb{E}_P\!\left[
 \int_t^\infty \beta_u \pi_u\,du \;\middle|\; \mathcal{A}_t
\right]
= \mathbb{E}_P\!\left[
 \int_0^\infty \beta_u \pi_u\,du \;\middle|\; \mathcal{A}_t
\right] - \int_0^t \beta_u \pi_u\,du.
\]
O primeiro termo c uma esperan SSa condicional (martingal), e o segundo
 c um processo de varia SS o finita. Portanto,
$\beta_t V_t^\pi + \int_0^t \beta_u \pi_u\,du$ c a soma de um
martingal com um termo de varia SS o finita que cancela exatamente a
parte de varia SS o finita da integral, resultando em um martingal.
\end{proof}

Esta proposi SS o fornece a equa SS o din mica que governa a evolu SS o do
deflator $D_t$ --- uma equa SS o que, no contexto da GAT, ser i
interpretada como a condi SS o de que a \textbf{conex o de mercado} tem
\textbf{curvatura zero}.

% ----------------------------------------------------------------------
\section{Resumo do Cap -tulo}
% ----------------------------------------------------------------------

O modelo cl issico de mercado em tempo cont -nuo assenta sobre seis pilares:

\begin{enumerate}
 \item \textbf{Espa SSo filtrado:} $(\Omega, \mathcal{A}, \mathbb{A}, P)$
 com condi SS ues usuais --- a base matem itica para a chegada
 gradual de informa SS o.
 \item \textbf{Pre SSos como semimartingales:} $S_t^j$ --- a classe
 maximal para a qual a integral estoc istica c bem definida.
 \item \textbf{Estrat cgias previs -veis e admiss -veis:} $x \in
 \mathcal{A}$ --- causalidade (n o-antecipa SS o) e limita SS o de
 perdas (exclus o de doubling strategies).
 \item \textbf{Autofinanciamento:} $dV_t = x_t \cdot dS_t$ (It ') ou
 $dV_t = x_t \circ dS_t$ (Stratonovich) --- a varia SS o da
 riqueza vem exclusivamente da varia SS o dos pre SSos.
 \item \textbf{NFLVR:} $\overline{C} \cap L_+^\infty = \{0\}$ --- a
 condi SS o topol 3gica que exclui n o apenas arbitragens, mas
 tamb cm suas aproxima SS ues.
 \item \textbf{FTAP:} NFLVR $\Leftrightarrow$ exist ancia de EMM
 $\Leftrightarrow$ exist ancia de state price deflator $\beta$
 --- o teorema que conecta efici ancia de mercado, medidas de
 probabilidade e apre SSamento.
\end{enumerate}

A partir do pr 3ximo cap -tulo, a GAT reformular i cada um destes pilares
na linguagem da geometria diferencial. A intui SS o essencial c: a
medida martingal equivalente $Q$ (ou o pricing kernel $\beta$) n o c
apenas uma ferramenta de apre SSamento --- c uma \textbf{se SS o de um
fibrado principal}, e a condi SS o NFLVR c a afirma SS o de que este
fibrado c \textbf{plano} (curvatura zero). O que a matem itica
financeira cl issica chama de ``mudan SSa de numer irio'' ou ``escolha de
conta caixa'', a GAT chama de \textbf{transforma SS o de gauge}.

\begin{center}
\fbox{\begin{minipage}{0.88\textwidth}
\textbf{Mensagem Central do Cap -tulo:} A teoria cl issica de arbitragem
 c completa e autocontida --- mas est i formulada em uma linguagem
(probabilidade) que \textit{esconde} a estrutura geom ctrica subjacente.
A GAT n o contradiz a teoria cl issica; ela a \textbf{revela} como um
caso particular ($R \equiv 0$) de uma estrutura muito mais rica, onde
arbitragem, holonomia e curvatura s o manifesta SS ues do mesmo fen 'meno
geom ctrico.
\end{minipage}}
\end{center}

% ======================================================================
% REFER SNCIAS DO CAP TULO (selecao principal)
% ======================================================================
\begin{thebibliography}{99}

\bibitem{delbaen1994}
Delbaen, F.; Schachermayer, W.
\textit{A General Version of the Fundamental Theorem of Asset Pricing}.
Mathematische Annalen, 300(1):463--520, 1994.
DOI: \texttt{10.1007/BF01450498}.

\bibitem{delbaen2006}
Delbaen, F.; Schachermayer, W.
\textit{The Mathematics of Arbitrage}.
Springer Finance, Berlin, 2006.
DOI: \texttt{10.1007/978-3-540-31299-4}.

\bibitem{dellacherie1972}
Dellacherie, C.
\textit{Capacit cs et Processus Stochastiques}.
Springer, Berlin, 1972.

\bibitem{farinelli2021}
Farinelli, S.
\textit{Geometric Arbitrage Theory and Market Dynamics Reloaded}.
arXiv:0910.1671v10, 2021.

\bibitem{harrison1979}
Harrison, J.M.; Kreps, D.M.
\textit{Martingales and Arbitrage in Multiperiod Securities Markets}.
Journal of Economic Theory, 20(3):381--408, 1979.
DOI: \texttt{10.1016/0022-0531(79)90043-7}.

\bibitem{harrison1981}
Harrison, J.M.; Pliska, S.R.
\textit{Martingales and Stochastic Integrals in the Theory of
Continuous Trading}.
Stochastic Processes and their Applications, 11(3):215--260, 1981.
DOI: \texttt{10.1016/0304-4149(81)90026-0}.

\bibitem{protter2005}
Protter, P.
\textit{Stochastic Integration and Differential Equations}.
Springer, 2nd ed., Berlin, 2005.
DOI: \texttt{10.1007/978-3-662-10061-5}.

\end{thebibliography}

% ======================================================================
% FIM DO CAPITULO 4
% ======================================================================


